% Содержимое подраздела 1.2 - Рекуррентные нейронные сети как генеративные модели "химического языка"


\section{Моделирование молекулярных структур как последовательностей с использованием рекуррентентных нейронных сетей}
\label{sec:rnn}

Исторически первыми и наиболее влиятельными архитектурами для генерации молекул стали рекуррентные нейронные сети (RNN), в частности их более сложная разновидность -- сети с долгой краткосрочной памятью (Long Short-Term Memory, LSTM). Принцип их работы в контексте химии основан на концепции авторегрессионной языковой модели. SMILES-строка разбивается на последовательность токенов (отдельных символов), и модель обучается на каждом шаге предсказывать следующий токен на основе всех предыдущих. В процессе генерации модель, начиная со стартового токена, посимвольно <<пишет>> новую SMILES-строку, постоянно опираясь на уже сгенерированную часть последовательности~\cite{wigh2022review}. Ранние работы в этой области продемонстрировали высокую эффективность, масштабируемость и, что самое главное, практическую применимость данного подхода.

Одним из первых масштабных исследований, доказавших состоятельность RNN-подхода, стала работа Ertl et al. (2017)~\cite{ertl2017silico}. Обучив относительно простую двухслойную LSTM-сеть на 509,000 биоактивных молекул из базы данных ChEMBL, авторы продемонстрировали возможность генерации \textbf{одного миллиона химически корректных молекул менее чем за два часа}. Несмотря на то, что только 32\% от всех сгенерированных строк оказались валидными, новизна результата была чрезвычайно высока: лишь 0.3\% сгенерированных структур совпадали с молекулами из обучающего набора. Ключевым достижением работы стала многоуровневая валидация, которая показала, что распределения физико-химических свойств (молекулярный вес, LogP, TPSA) и оценки синтетической доступности (SA score) у сгенерированных молекул практически идентичны реальным соединениям из ChEMBL. Наиболее убедительным доказательством стало применение QSAR-моделей, которое показало, что новые молекулы имеют сопоставимый с известными лекарствами потенциал биологической активности.

В то время как работа Ertl et al. продемонстрировала масштабируемость, исследование Bjerrum и Threllfall (2017)~\cite{bjerrum2017molecular} было сфокусировано на детальном анализе качества и новизны генерируемых структур. Обучая модель на различных подмножествах базы ZINC (<<фрагменто-подобные>> и <<лекарственно-подобные>> молекулы), они показали, что RNN способна точно воспроизводить распределения свойств из обучающей выборки. Это доказывает, что модель улавливает скрытые <<химические правила>>, заложенные в данных. Анализ новизны показал, что до \textbf{83\%} сгенерированных молекул отсутствовали в исходном обучающем наборе, при этом их оценка синтетической доступности (SA score) имела такое же распределение, как и у реальных соединений. Таким образом, было доказано, что модель генерирует не просто валидные, но и химически осмысленные, новые и потенциально синтезируемые структуры.

Кульминацией ранних исследований стала работа, продемонстрировавшая переход от теоретических выкладок к реальной экспериментальной проверке~\cite{merk2018novo}. В этом исследовании был применен подход на основе переноса обучения (transfer learning): RNN-модель, предварительно обученная на большой базе данных ChEMBL, была дообучена на очень малом наборе из 25 молекул-агонистов ядерных рецепторов RXR и PPAR. Из сгенерированных кандидатов пять были отобраны для химического синтеза. В результате \textbf{четыре из пяти соединений продемонстрировали желаемую биологическую активность} в наномолярном и микромолярном диапазоне, причем одно из соединений показало активность \texttt{EC\textsubscript{50} = 0.06 $\mu$M}. Эта работа стала первым проспективным применением генеративной модели, завершившим полный цикл от дизайна с помощью ИИ до синтеза и успешной биологической валидации, и доказала практическую ценность RNN-подхода для de novo дизайна лекарств.
